\DPrinciple{原则七：坚持开源普惠}{“智能为公，成果共享”}
推动开源技术广泛应用于政府治理、公共卫生、教育公平等领域，服务社会治理现代化与民生改善。鼓励开源成为数字公共品的载体，建立以贡献为导向的激励机制，让创新成果惠及所有人，尤其是边缘群体。

\DHead{IV. 行动共识：迈向全球 AI 治理新格局}
我们呼吁：各国政府、科技企业、学界与公民社会携手同行，以本宣言为基础，深化在 AI 治理、算力共享、人才培育等领域的务实合作，争取早期收获。

我们支持依托联合国及相关国际机制，围绕人工智能伦理准则、跨境数据流动规则、AI 版权保护等紧迫议题展开对话，凝聚共识，推动开源成为构建公正合理全球人工智能治理体系的生动实践。

\DHead{结语}
开源，是我们对开放精神的坚守，是对人类团结的信念，是对更美好世界的共同追求。让我们以协作代替割裂，以透明回应质疑，以共享消弭垄断。共同构建一个更加开放、公平、可持续的智能文明。

开源创新，属于全人类；智能未来，由我们共同创造。
